% Corrected French for printed p.195.
\clearpage
\SourcePage{195}

\LayerTarget{tohoku:proof:remark-4.4.1:1}
Comme $\mathbf{Ext}_{\mathbf Z(G)}^n(\mathbf Z,A)$ est le faisceau associé au
préfaisceau des $\mathrm{Ext}_{\mathbf Z(G)}^n(U;\mathbf Z,A)$ qui est nul
d'après ce qui précède (puisque $A|U$ est aussi injectif), notre assertion est
démontrée.

\LayerTarget{tohoku:remark:4.4.1:2}
2. Supposons que $A$ soit le faisceau \textit{constant} défini par un
$G$-module $M$, alors on écrira évidemment $H^p(X;G,M)$ et
$\mathbf H^p(G,M)$ au lieu de $H^p(X;G,A)$ et $\mathbf H^p(G,A)$. Si $G$
opère trivialement sur $M$, alors $H^p(X;G,M)$ se calcule donc complètement
par la formule 4.4.3, ou la formule 4.4.2 si $M$ est un espace vectoriel sur
un corps $k$. Dans ce dernier cas, et si $X$ est un espace HLC, on obtient
même :
\LayerTarget{tohoku:formula:bilinear-equivariant-cohomology}
$H^*(X,G,M)=$ espace des applications bilinéaires de
$H_*(G,k)\times H_*(X,k)$ dans $M$.
On fera attention qu'en général, $\mathbf H^p(G,M)$ n'est pas le faisceau
constant défini par $H^p(G,M)$, cf. remarque 2 de 4.3.

\LayerTarget{tohoku:chapter:V}
\bigskip
\begin{center}
{\large\bfseries Chapitre V\quad Etude cohomologique des espaces à
opérateurs.}
\end{center}

\LayerTarget{tohoku:section:5.1}
\textbf{5.1. Généralités sur les $G$-faisceaux.}
\LayerTarget{tohoku:setup:g-space-and-orbit-space}
Dans tout ce Chapitre, nous considérons un espace $X$ où opère un groupe $G$
(à gauche pour fixer les idées). L'opération définie par un $g\in G$ sera
notée $x\to g\mathbin{\cdot}x$. Nous n'exigeons pas que $G$ opère
fidèlement ; en particulier, nous aurons à considérer aussi le cas où $G$
opérerait trivialement. Nous écrirons $X(G)$ pour $X$ muni de la structure
supplémentaire définie par les opérations de $G$. L'espace des trajectoires
$X/G$ sera noté $Y$, il sera muni de la topologie quotient. L'application
canonique de $X$ sur $Y$ sera notée $f$, c'est une application continue
ouverte. $Y$ sera considéré comme muni du groupe d'opérateurs $G$ opérant
trivialement, on écrira donc $Y(G)$ pour rappeler cette structure
supplémentaire sur $Y$.

\LayerTarget{tohoku:definition:g-sheaf}
Nous appellerons \textit{$G$-faisceau} sur $X=X(G)$ un faisceau (d'ensembles)
$A$ sur $X$, dans lequel $G$ opère de façon compatible avec ses opérations
sur $X$. Pour donner un sens à cette définition, on pourra par exemple
considérer $A$ comme espace étalé dans $X$ (cf. 3.1) ; nous n'insisterons pas.
On définit de même la notion de \textit{$G$-faisceau de groupes, d'anneaux,
de $G$-faisceau abélien} (= $G$-faisceau de groupes abéliens) etc.
\LayerTarget{tohoku:example:g-sheaves}
De façon imagée, on peut dire que si $X$ est muni de certaines structures et
si un faisceau $A$ sur $X$ est défini en termes structuraux, alors $A$ est un
$G$-faisceau de façon naturelle si les opérations de $G$ dans $X$ sont des
automorphismes. Ainsi, un faisceau constant peut toujours être considéré comme
un $G$-faisceau (``\textit{$G$-faisceau trivial}''), de même le faisceau des germes
d'applications quelconques resp. continues de $X$ dans un ensemble donné, le
faisceau des germes de fonctions holomorphes si $X$ est une variété holomorphe
et si les opérations de $G$ sont des automorphismes de la variété $X$, etc.
\LayerTarget{tohoku:definition:g-homomorphism}
On appelle \textit{$G$-homomorphisme} d'un faisceau dans un autre un
homomorphisme de faisceau qui permute aux opérations de $G$. Si les faisceaux
envisagés sont des faisceaux de groupes par exemple, on sous-entend que
l'homomorphisme respecte cette structure, comme d'habitude.
\LayerTarget{tohoku:definition:category-g-sheaves}
Avec cette notion d'homomorphisme, les $G$-faisceaux d'ensembles (resp. les
$G$-faisceaux de groupes etc.) forment une \textit{catégorie} (cf. 1.1), qui a
les mêmes propriétés que la catégorie correspondante sans groupe d'opérateurs
$G$, (qui en est d'ailleurs un cas particulier, correspondant au cas où $G$
est réduit à l'élément neutre).

\TohokuContinuation{196}
